% Part: model-theory
% Chapter: interpolation
% Section: interpolation-proof

\documentclass[../../../include/open-logic-section]{subfiles}

\begin{document}

% SEGMENT OLP-0201-S01

\olfileid{mod}{int}{prf}

\olsection{கிரெய்கின் இடையீட்டுத் தேற்றம்}

\begin{thm}[கிரெய்கின் இடையீட்டுத் தேற்றம்]
\ollabel{thm:interpol}
$\Entails !A \lif !B$ எனில், $\Entails !A \lif !C$ மற்றும்
$\Entails !C \lif !B$ ஆகுமாறு !!a{sentence} $!C$ ஒன்று உள்ளது;
மேலும் $!C$ இலுள்ள ஒவ்வொரு !!{constant}, !!{function},
!!{predicate} உம் ($\eq$ தவிர) $!A$ மற்றும்~$!B$ ஆகிய இரண்டிலும்
இடம்பெறும். !!{sentence} $!C$ ஆனது $!A$ மற்றும்~$!B$ இன்
\emph{இடையீட்டி} எனப்படுகிறது.
\end{thm}

\begin{proof}
$!A$ இன் மொழி $\Lang{L}_1$ என்றும், $!B$ இன் மொழி $\Lang{L}_2$
என்றும் கொள்க. $\Lang{L}_0 = \Lang{L}_1 \cap \Lang{L}_2$ என்க.
ஒவ்வொரு $i \in \{0, 1, 2 \}$ க்கும், எண்ணிலடங்கா புதிய
!!{constant}s $\Obj c_0,
\Obj c_1, \Obj c_2, \dots$ ஐ $\Lang{L}_i$ உடன் சேர்ப்பதன் மூலம்
$\Lang{L}'_i$ பெறப்படட்டும்.

$!A$ நிறைவுறத்தகாதது எனில், $\lexists[x][\eq/[x][x]]$ ஓர் இடையீட்டி.
$\lnot !B$ நிறைவுறத்தகாதது எனில் (எனவே $!B$ செல்லுபடியானது எனில்),
$\lexists[x][\eq[x][x]]$ ஓர் இடையீட்டி. ஆகவே $!A$ மற்றும்
$\lnot !B$ ஆகிய இரண்டும் நிறைவுறத்தக்கவை என்றும் கொள்ளலாம்.

இடையீட்டுத் தேற்றத்தின் முரண்நிலையை நிரூபிக்க, $!A$ மற்றும் $!B$
க்கான இடையீட்டி இல்லை என்று கொள்க. வேறுவிதமாகச் சொன்னால்,
$\{!A\}$ மற்றும் $\{\lnot !B\}$ ஆகியவை $\Lang{L}_0$ இல்
பிரிக்கவியலாதவை என்று கொள்க.

சோடி $(\{ !A \}, \{\lnot!B\})$ ஐ மீப்பெருமப் பிரிக்கவியலாச் சோடி
$(\Gamma^*, \Delta^*)$ ஆக நீட்டிப்பதே நமது இலக்கு. $\Lang{L}_1$
இன் !!{sentence}s ஐ $!A_0$, $!A_1$, $!A_2$, \dots உம்,
~$\Lang{L}_2$ இன் !!{sentence}s ஐ $!B_0$, $!B_1$, $!B_2$, \dots உம்
வரிசைப்படுத்தட்டும். !!{sentence}s இன் கணங்களின் இரு ஏறு தொடர்கள்
$(\Gamma_n, \Delta_n)$ ஐ $n \ge 0$ க்கு பின்வருமாறு வரையறுக்கிறோம்.
$\Gamma_0 = \{ !A\}$ மற்றும் $\Delta_0 = \{\lnot !B \}$ என்க.
$(\Gamma_n, \Delta_n)$ ஏற்கெனவே வரையறுக்கப்பட்டதாகக் கொண்டு,
$\Gamma_{n+1}$ மற்றும் $\Delta_{n+1}$ ஐப் பின்வருமாறு வரையறுக்க:
\begin{enumerate}
\item $\Gamma_n \cup \{!A_n \}$ மற்றும் $\Delta_n$ ஆகியவை
  $\Lang{L}'_0$ இல் பிரிக்கவியலாதவை எனில், $!A_n$ ஐ
  $\Gamma_{n+1}$ இல் இடுக. மேலும் $!A_n$ ஆனது இருத்தலளவையீட்டு
  !!{formula} $\lexists[x][!S]$ எனில், $\Gamma_n$, $\Delta_n$,
  $!A_n$, $!B_n$ ஆகியவற்றில் இடம்பெறாத புதிய !!{constant} $c$
  ஒன்றைத் தேர்ந்து, $\Subst{!S}{c}{x}$ ஐ $\Gamma_{n+1}$ இல் இடுக.
\item $\Gamma_{n+1}$ மற்றும் $\Delta_n \cup \{!B_n \}$ ஆகியவை
  $\Lang{L}'_0$ இல் பிரிக்கவியலாதவை எனில், $!B_n$ ஐ
  $\Delta_{n+1}$ இல் இடுக. மேலும் $!B_n$ ஆனது இருத்தலளவையீட்டு
  !!{formula} $\lexists[x][!S]$ எனில், $\Gamma_{n+1}$, $\Delta_n$,
  $!A_n$, $!B_n$ ஆகியவற்றில் இடம்பெறாத புதிய !!{constant} $c$
  ஒன்றைத் தேர்ந்து, $\Subst{!S}{c}{x}$ ஐ $\Delta_{n+1}$ இல் இடுக.
\end{enumerate}
இறுதியாக, பின்வருமாறு வரையறுக்க:
\begin{align*}
  \Gamma^* & = \bigcup_{n\ge 0} \Gamma_n, &
  \Delta^* & = \bigcup_{n\ge 0} \Delta_n.
\end{align*}
$n$ மீதான ஒருங்குநிகழ் தொகுத்தறிதலால் இப்போது பின்வருவனவற்றை
நிரூபிக்கலாம்:
\begin{enumerate}
\item\ollabel{part-a} $\Gamma_n$ மற்றும் $\Delta_n$ ஆகியவை
  $\Lang{L}'_0$ இல் பிரிக்கவியலாதவை;
\item\ollabel{part-b} $\Gamma_{n+1}$ மற்றும் $\Delta_n$ ஆகியவை
  $\Lang{L}'_0$ இல் பிரிக்கவியலாதவை.
\end{enumerate}
\olref{part-a} இன் அடிப்படை நிலையை \olref[sep]{lem:sep1} தருகிறது.
\olref{part-b} பகுதிக்கு மூன்று நிலைகளை வேறுபடுத்த வேண்டும்:
\begin{enumerate}
\item $\Gamma_0 \cup \{!A_0 \}$ மற்றும் $\Delta_0$ பிரிக்கத்தக்கவை
  எனில், $\Gamma_1 = \Gamma_0$; அப்போது \olref{part-b} என்பது
  \olref{part-a} மட்டுமே;
\item $\Gamma_1 = \Gamma_0 \cup\{ !A_0\}$ எனில், கட்டுமானத்தின்படி
  $\Gamma_1$ மற்றும் $\Delta_0$ பிரிக்கவியலாதவை.
\item $!A_0$ இருத்தலளவையீட்டுடையதாக இருந்து,
  $\Gamma_1 = \Gamma_0 \cup \{ \lexists[x][!S], \Subst{!S}{c}{x}
  \}$ ஆகும் நிலையை இன்னும் கருத வேண்டும். கட்டுமானத்தின்படி,
  $\Gamma_0 \cup \{ \lexists[x][!S]\}$ மற்றும் $\Delta_0$
  பிரிக்கவியலாதவை; ஆகவே \olref[sep]{lem:sep2} படி,
  $\Gamma_0 \cup \{ \lexists[x][!S], \Subst{!S}{c}{x} \}$ மற்றும்
  $\Delta_0$ உம் பிரிக்கவியலாதவை.
\end{enumerate}
இதனால் மேலுள்ள \olref{part-a} மற்றும் \olref{part-b} க்கான
தொகுத்தறிதலின் அடிப்படை நிலை நிறைவடைகிறது. இப்போது தொகுத்தறிதல் படி.
\olref{part-a} க்கு, $\Delta_{n+1} = \Delta_n \cup \{ !B_n \}$
எனில், $\Gamma_{n+1}$ மற்றும் $\Delta_{n+1}$ கட்டுமானத்தின்படி
பிரிக்கவியலாதவை ($!B_n$ இருத்தலளவையீட்டுடையதாக இருந்தாலும்கூட
\olref[sep]{lem:sep2} படி); $\Delta_{n+1} = \Delta_n$ எனில்
($\Gamma_{n+1}$ மற்றும் $\Delta_n \cup \{!B_n\}$ பிரிக்கத்தக்கவை
என்பதால்), \olref{part-b} மீதான தொகுத்தறிதல் கருதுகோளைப்
பயன்படுத்துகிறோம். \olref{part-b} க்கான தொகுத்தறிதல் படியில்,
$\Gamma_{n+2} = \Gamma_{n+1} \cup
\{!A_{n+1} \}$ எனில், $\Gamma_{n+2}$ மற்றும் $\Delta_{n+1}$
கட்டுமானத்தின்படி பிரிக்கவியலாதவை ($!A_{n+1}$
இருத்தலளவையீட்டுடையதாக இருந்தாலும்கூட \olref[sep]{lem:sep2} படி);
$\Gamma_{n+2} = \Gamma_{n+1}$ எனில், இப்போது நிரூபித்த
\olref{part-a} க்கான தொகுத்தறிதல் நிலையைப் பயன்படுத்துகிறோம்.
இதனால் \olref{part-a} மற்றும் \olref{part-b} மீதான தொகுத்தறிதல்
நிறைவடைகிறது.

இதிலிருந்து $\Gamma^*$ மற்றும் $\Delta^*$ பிரிக்கவியலாதவை என்று
பின்வருகிறது; இல்லையெனில் இறுதிக்குட்பட்ட தன்மையின்படி,
$\Gamma_n$ மற்றும் $\Delta_n$ ஐப் பிரிக்கும் $n \ge 0$ ஒன்று
இருக்கும்; இது \olref{part-a} க்கு மாறானது. குறிப்பாக,
$\Gamma^*$ மற்றும் $\Delta^*$ முரணின்மையுடையவை: முதலாவதோ
பிந்தையதோ முரண்பட்டதாக இருந்தால், முறையே
$\lexists[x][\eq/[x][x]]$ அல்லது
$\lforall[x][\eq[x][x]]$ அவற்றைப் பிரிக்கும்.

இப்போது $\Gamma^*$ ஆனது $\Lang{L}'_1$ இல் மீப்பெரும
முரணின்மையுடையது என்றும், அவ்வாறே $\Delta^*$ ஆனது
$\Lang{L}'_2$ இல் மீப்பெரும முரணின்மையுடையது என்றும் காட்டுவோம்.
முதலாவதற்கு, ஏதேனும் $n \ge 0$ க்கு $!A_n \notin \Gamma^*$ மற்றும்
$\lnot !A_n
\notin \Gamma^*$ என்று கொள்க. $!A_n \notin \Gamma^*$ எனில்,
$\Gamma_n \cup \{!A_n \}$ ஆனது $\Delta_n$ இலிருந்து பிரிக்கத்தக்கது;
எனவே பின்வரும் இரண்டும் பொருந்துமாறு $!C \in \Lang{L}'_0$ ஒன்று உண்டு:
\begin{align*}
  \Gamma^* & \Entails !A_n \lif !C, &
  \Delta^* & \Entails \lnot !C.
\end{align*}
அவ்வாறே, $\lnot !A_n \notin \Gamma^*$ எனில், பின்வரும் இரண்டும்
பொருந்துமாறு $!C' \in
\Lang{L}'_0$ ஒன்று உண்டு:
\begin{align*}
  \Gamma^* & \Entails \lnot !A_n \lif !C', &
  \Delta^* & \Entails \lnot !C'.
\end{align*}
கூற்றுக்கணிப்புத் தருக்கப்படி, $\Gamma^* \Entails !C \lor !C'$ மற்றும்
$\Delta^* \Entails \lnot (!C \lor !C')$; எனவே $!C \lor
!C'$ ஆனது $\Gamma^*$ மற்றும் $\Delta^*$ ஐப் பிரிக்கிறது.
இதேபோன்ற வாதம் $\Delta^*$ மீப்பெருமமானது என்று நிறுவுகிறது.

இறுதியாக, $\Gamma^* \cap \Delta^*$ ஆனது $\Lang{L}'_0$ இல்
மீப்பெரும முரணின்மையுடையது என்று காட்டுவோம். அது முரணின்மையுடைய
கணங்களின் வெட்டு என்பதால் வெளிப்படையாக முரணின்மையுடையது.
மீப்பெருமத்தை காட்ட, $!S \in
\Lang{L}'_0$ என்க. இப்போது $\Gamma^*$ ஆனது $\Lang{L'_1}
\supseteq \Lang{L'_0}$ இல் மீப்பெருமமானது; அவ்வாறே $\Delta^*$ ஆனது
$\Lang{L'_2} \supseteq \Lang{L'_0}$ இல் மீப்பெருமமானது. எனவே
$!S \in \Gamma^*$ அல்லது $\lnot !S \in \Gamma^*$; மேலும்
$!S \in \Delta^*$ அல்லது $\lnot !S \in \Delta^*$. $!S \in
\Gamma^*$ மற்றும் $\lnot !S \in \Delta^*$ எனில், $!S$ ஆனது
$\Gamma^*$ மற்றும் $\Delta^*$ ஐப் பிரிக்கும்; $\lnot !S \in
\Gamma^*$ மற்றும் $!S \in \Delta^*$ எனில், $\lnot !S$ ஆனது
$\Gamma^*$ மற்றும் $\Delta^*$ ஐப் பிரிக்கும். எனவே $!S \in
\Gamma^* \cap \Delta^*$ அல்லது $\lnot !S \in \Gamma^* \cap \Delta^*$;
ஆகவே $\Gamma^* \cap \Delta^*$ மீப்பெருமமானது.

$\Gamma^*$ மீப்பெரும முரணின்மையுடையது என்பதால், நிறைவுத் தேற்றத்தின்
நிரூபிப்பைப் போலவே (\olref[fol][com][cth]{thm:completeness}),
!!{constant}s ஐப் பொருள்கொள்ளும் உறுப்புகள் $\Assign{c}{M'_1}$
அனைத்தையும் அவற்றை மட்டுமே கொண்ட !!{domain} $\Domain{M'_1}$ உடைய
மாதிரி $\Struct{M}'_1$ அதற்கு உண்டு. அவ்வாறே, $\Delta^*$ க்கு,
!!{constant}s இன் பொருள்கோள்கள் $\Assign{c}{M'_2}$ ஆல் தரப்படும்
!!{domain} $\Domain{M'_2}$ உடைய மாதிரி $\Struct{M}'_2$ உள்ளது.

$\Lang{L'_1} \setminus \Lang{L'_0}$ இலுள்ள !!{constant}s,
!!{function}s, !!{predicate}s ஆகியவற்றின் பொருள்கோள்களை நீக்குவதன்
மூலம் $\Struct{M'_1}$ இலிருந்து $\Struct{M_1}$ பெறப்படட்டும்;
$\Struct{M_2}$ உம் அவ்வாறே. அப்போது
$h(\Assign{c}{M'_1}) = \Assign{c}{M'_2}$ என்று வரையறுக்கப்படும்
$h \colon M_1 \to M_2$ ஆனது $\Lang{L}'_0$ இல் ஓர் ஓரினச்சார்பு;
ஏனெனில் காட்டியபடி $\Gamma^* \cap \Delta^*$ ஆனது $\Lang{L}'_0$
இல் மீப்பெரும முரணின்மையுடையது. எந்த $\Lang{L}'_0$-!!{sentence} உம்
$\Gamma^*$ மற்றும் $\Delta^*$ ஆகிய இரண்டிலும் இடம்பெறும் அல்லது
இரண்டிலும் இடம்பெறாது என்பதால் இது பின்வருகிறது: எனவே
$\Assign{c}{M'_1} \in
\Assign{P}{M'_1}$ எனில், எனில் மட்டுமே, $\Atom{P}{c} \in \Gamma^*$;
எனில், எனில் மட்டுமே, $\Atom{P}{c} \in \Delta^*$; எனில், எனில்
மட்டுமே, $\Assign{c}{M'_2} \in
\Assign{P}{M'_2}$. ஓரினச்சார்புகள் நிறைவுறுத்தும் பிற நிபந்தனைகளையும்
இதேபோல் நிறுவலாம்.

இப்போது !!{language} $\Lang{L_1} \cup \Lang{L_2}$ க்கான மாதிரி
$\Struct{M}$ ஒன்றைப் பின்வருமாறு வரையறுப்போம்:
\begin{enumerate}
\item !!{domain} $\Domain{M}$ என்பது $\Domain{M_2}$ தான்; அதாவது
  எல்லா உறுப்புகள் $\Assign{c}{M'_2}$ இன் கணம்;
\item !!a{predicate}~$P$ ஆனது $\Lang{L_2} \setminus
  \Lang{L_1}$ இல் இருந்தால், $\Assign{P}{M} = \Assign{P}{M'_2}$;
\item பயனிலை $P$ ஆனது $\Lang{L}_1\setminus \Lang{L}_2$ இல் இருந்தால்,
  $\Assign{P}{M} = h(\Assign{P}{M'_2})$; அதாவது,
  $\tuple{\Assign{c_1}{M'_2}, \dots, \Assign{c_n}{M'_2}} \in
  \Assign{P}{M}$ எனில், எனில் மட்டுமே,
  $\tuple{\Assign{c_1}{M'_1}, \dots,
  \Assign{c_n}{M'_1}} \in \Assign{P}{M'_1}$.
\item !!a{predicate} $P$ ஆனது $\Lang{L}_0$ இல் இருந்தால்,
  $\Assign{P}{M} =
  \Assign{P}{M'_2} = h(\Assign{P}{M'_1})$.
\item !!^{function}s மற்றும் !!{constant}s உட்பட,
  $\Lang{L}_1 \cup \Lang{L}_2$ இன் சார்புகள் இதேபோல் கையாளப்படுகின்றன.
\end{enumerate}

இறுதியாக, $\Struct{M}$ ஆனது $\Lang{L'_1}$ இன் எல்லா !!{formula}s
மீதும் $\Struct{M'_1}$ உடனும், $\Lang{L'_2}$ இன் எல்லா !!{formula}s
மீதும் $\Struct{M'_2}$ உடனும் உடன்படுகிறது என்பதை !!{formula}s மீது
தொகுத்தறிந்து காட்டலாம். குறிப்பாக,
$\Struct{M} \Entails \Gamma^* \cup \Delta^*$; எனவே $\Struct{M}
\Entails !A$ மற்றும் $\Struct{M} \Entails \lnot!B$; ஆகவே
$\not\Entails !A \lif !B$. இதனால் கிரெய்கின் இடையீட்டுத் தேற்றத்தின்
நிரூபிப்பு நிறைவடைகிறது.
\end{proof}

\end{document}
